\documentclass[11pt]{letter}
\usepackage[a4paper,margin=1in]{geometry}
\usepackage{setspace}
\setstretch{1.05}

\begin{document}

\begin{letter}{Editor-in-Chief\\
Applied Physics Reviews\\
AIP Publishing}

\opening{Dear Editor-in-Chief,}

I respectfully submit the manuscript entitled
\textbf{``Orientation-Dependent Dispersion and Wave Steering in Two-Dimensional Strain-Gradient Elasticity: Directional Spectral Separation versus Complete Bragg Gaps under Material Contrast''}
as an Original Research Article for consideration in \textit{Applied Physics Reviews}.

\textbf{Knowledge gap and novelty.} The work addresses the need for a carefully verified framework that can distinguish orientation-driven dispersion and energy-flow steering in a homogeneous higher-order elastic medium from complete Bragg attenuation caused by periodic material contrast. I combine a rotated ellipsoidal characteristic-length tensor, local Mindlin Form-II strain-gradient elasticity, independent gradient micro-inertia, and a conforming 32-degree-of-freedom Bogner--Fox--Schmit element. Bloch phases are imposed on displacement and derivative degrees of freedom, and the study evaluates both directional spectral and complete-gap measures rather than treating them as interchangeable.

\textbf{Importance of the results.} In the homogeneous case, the largest sampled $\Gamma$--$X$ directional spectral separation is 0.0677, while the complete-gap measure remains negative; the maximum acoustic-branch group-velocity deflection is 2.7846$^\circ$. For the epoxy--YBCO contrast case, a positive complete gap is obtained, but its reported width decreases from 2.5732 on a $4\times4$ mesh to 1.9179 on a $64\times64$ mesh. I explicitly identify this result as mesh-unconverged. The distinction clarifies which response is attributable to tensor orientation and which requires material contrast, while the numerical checks and disclosed resolution limits help readers assess the evidence and its scope.

\textbf{Appeal to APR readers.} The framework and findings connect applied continuum physics, computational mechanics, phononic materials, and elastic-wave design. I believe the separation of wave steering from complete attenuation will be of interest to APR readers working on engineered media and wave control.

\textbf{Related work by the author.} I confirm that there is no closely related published or submitted work by me to list.

\textbf{Submission declarations.} I confirm that this manuscript is original, has not been published previously, and is not under consideration elsewhere. I approve its submission to Applied Physics Reviews.

Thank you for considering my manuscript.

\closing{Sincerely,}

\textbf{Vipin Gupta}\\
Corresponding Author\\
Department of Mathematics\\
Gurugram University, Gurugram-122505, India\\
E-mail: vipin@gurugramuniversity.ac.in

\end{letter}

\end{document}
